% ---
% Capa
% ---
\imprimircapa
% ---

% ---
% Folha de rosto
% (o * indica que haverá a ficha bibliográfica)
% ---
\imprimirfolhaderosto*
% ---

% ---
% Inserir a ficha bibliografica
% ---
% http://ficha.bu.ufsc.br/
\begin{fichacatalografica}
  \includepdf{beforetext/Ficha_Catalografica.pdf}
\end{fichacatalografica}
% ---

% ---
% Inserir folha de aprovação
% ---
\begin{folhadeaprovacao}
  \OnehalfSpacing
  \centering
  \imprimirautor\\%
  \vspace*{10pt}    
  \textbf{\imprimirtitulo}%
  \ifnotempty{\imprimirsubtitulo}{:~\imprimirsubtitulo}\\%
  %    \vspace*{31.5pt}%3\baselineskip
  \vspace*{\baselineskip}
  %\begin{minipage}{\textwidth}
  % ~do~\imprimirprograma~do~\imprimircentro~da~\imprimirinstituicao~para~a~obtenção~do~título~de~\imprimirformacao.
  Este~\imprimirtipotrabalho~foi julgado adequado para obtenção do Título de “\imprimirformacao” e aprovado em sua forma final pelo~\imprimirprograma. \\
    \vspace*{\baselineskip}
  \imprimirlocal, \imprimirdata. \\
  \vspace*{2\baselineskip}
  \assinatura{\OnehalfSpacing\imprimircoordenador \\ \imprimircoordenadorRotulo~do Curso}
  \vspace*{2\baselineskip}
  \textbf{Banca Examinadora:} \\
  \vspace*{\baselineskip}
  \assinatura{\OnehalfSpacing\imprimirorientador \\ \imprimirorientadorRotulo}
  \assinatura{\OnehalfSpacing\imprimircoorientador \\ \imprimircoorientadorRotulo}
  %\end{minipage}%
  \vspace*{\baselineskip}
  \assinatura{Me. Eduardo Willwock Lussi,\\
  Avaliador \\
  Universidade Federal de Santa Catarina}

  \vspace*{\baselineskip}
  \assinatura{Ma. Letícia Bertuzzi\\
  Avaliadora \\
  Universidade Federal de Santa Catarina}
\end{folhadeaprovacao}
% ---

% ---
% Dedicatória
% ---
%\begin{dedicatoria}
%  \vspace*{\fill}
%  \noindent
%  \begin{adjustwidth*}{}{5.5cm}     
%    Este trabalho é dedicado aos meus colegas de classe e aos meus queridos pais.
%  \end{adjustwidth*}
%\end{dedicatoria}
% ---

% ---
% Agradecimentos
%
%\begin{agradecimentos}
%  Inserir os agradecimentos aos colaboradores à execução do trabalho. 
%  
%  Xxxxxxxxxxxxxxxxxxxxxxxxxxxxxxxxxxxxxxxxxxxxxxxxxxxxxxxxxxxxxxxxxxxxxx. 
%\end{agradecimentos}
% ---

% ---
% Epígrafe
% ---
%\begin{epigrafe}
%  \vspace*{\fill}
%  \begin{flushright}
%    \textit{``Texto da Epígrafe.\\
%      Citação relativa ao tema do trabalho.\\
%      É opcional. A epígrafe pode também aparecer\\
%      na abertura de cada seção ou capítulo.\\
%      Deve ser elaborada de acordo com a NBR 10520.''\\
%      (Autor da epígrafe, ano)}
%  \end{flushright}
%\end{epigrafe}
% ---

% ---
% RESUMOS
% ---

% resumo em português

\setlength{\absparsep}{18pt} % ajusta o espaçamento dos parágrafos do resumo
\begin{resumo}
  Atualmente, o desenvolvimento de algoritmos quânticos é limitado pela falta de formas eficazes e acessíveis de executá-los na prática. Como o hardware quântico disponível ainda apresenta baixa capacidade de processamento e custo elevado de acesso, a simulação em hardware clássico é uma alternativa essencial para viabilizar o projeto, o teste e a depuração de algoritmos quânticos. Nesse contexto, destaca-se o Ket, uma plataforma de desenvolvimento quântico que permite aos usuários elaborar e simular seus programas localmente.

  Por outro lado, a simulação quântica clássica impõe desafios significativos: em abordagens convencionais, o consumo de memória cresce exponencialmente com o número de qubits do circuito, tornando a execução inviável para sistemas maiores. Para contornar essa limitação, este trabalho teve como objetivo implementar um método de simulação baseado na decomposição matricial das portas lógicas e integrá-lo à plataforma Ket. A estratégia faz uso do QMDD (Quantum Multiple-valued Decision Diagram), uma estrutura de dados capaz de compactar matrizes e vetores com alta redundância de valores. Essa característica tende a ser encontrada em circuitos com baixa taxa de entrelaçamento e permite que simulações sejam feitas com complexidade de espaço e tempo potencialmente subexponenciais.

  O simulador foi implementado na linguagem de programação Rust e integrado à plataforma Ket. Seus resultados foram validados ao serem comparados com os resultados de outros simuladores já presentes na plataforma ao executar algoritmos comuns da computação quântica. Além disso, o desempenho do simulador foi medido e comparado as alternativas já disponíveis, apresentando vantagem em circuitos com poucas portas controladas. Então, o objetivo principal do trabalho foi atingido: um novo simulador está disponível para usuários do Ket.

  \textbf{Palavras-chave}: Computação Quântica. Simulação Quântica. Diagrama de Decisão. Ket.
\end{resumo}

% resumo em inglês
\begin{resumo}[Abstract]
  \begin{otherlanguage*}{english}
    Currently, the development of quantum algorithms is limited by the of efficient and accessible means to execute them. As quantum hardware available still presents low processing capacity and high access cost, simulation via classic hardware is an essential alternative to make viable the design, testing and debugging of quantum algorithms. Hence, there is Ket, a quantum development plataform which allows users to create and simulate quantum programs locally.

    Similarly, quantum simulation in classic computers has significant challenges: in conventional approaches, memory usage increases exponentially with the number of qubits in the circuit, rendering it infeasible for larger circuits. To work around that limitation, this work has the goal of implementing a simulation method based on matrix decomposition of quantum gates and integrate it to Ket. This approach makes use of QMDD (Quantm Multiple-valued Decision Diagram), a data structure capable of compressing matrices and vectors with high redundancy of values. This characteristic tends to appear in circuits with low levels of entanglement and allows for simulations to be done with subexponential time and space complexity.

    The simulator was implemented in the Rust programming language and integrated to the Ket platform. Its results were validated by being compared to the ones of other simulators already available in the platform when executing common quantum algorithms. Furthermore, the simulator's performance was measured and compared with the existing alternatives, showing advantage in circuits with few controled gates. Therefore, the main goal of the work was reached: a new and efficient simulator is available to Ket users.

    \textbf{Keywords}: Quantum Computing. Quantum Simulation. Decision Diagram. Ket.
  \end{otherlanguage*}
\end{resumo}

%% resumo em francês 
%\begin{resumo}[Résumé]
% \begin{otherlanguage*}{french}
%    Il s'agit d'un résumé en français.
% 
%   \textbf{Mots-clés}: latex. abntex. publication de textes.
% \end{otherlanguage*}
%\end{resumo}
%
%% resumo em espanhol
%\begin{resumo}[Resumen]
% \begin{otherlanguage*}{spanish}
%   Este es el resumen en español.
%  
%   \textbf{Palabras clave}: latex. abntex. publicación de textos.
% \end{otherlanguage*}
%\end{resumo}
%% ---

{%hidelinks
  \hypersetup{hidelinks}
  % ---
  % inserir lista de ilustrações
  % ---
  \pdfbookmark[0]{\listfigurename}{lof}
  \listoffigures*
  \cleardoublepage
  % ---
  
  % ---
  % inserir lista de quadros
  % ---
  % \pdfbookmark[0]{\listofquadrosname}{loq}
  % \listofquadros*
  % \cleardoublepage
  % ---
  
  % ---
  % inserir lista de tabelas
  % ---
  % \pdfbookmark[0]{\listtablename}{lot}
  % \listoftables*
  % \cleardoublepage
  % ---
  
  % ---
  % inserir lista de abreviaturas e siglas (devem ser declarados no preambulo)
  % ---
  \imprimirlistadesiglas
  % ---
  
  % ---
  % inserir lista de símbolos (devem ser declarados no preambulo)
  % ---
  \imprimirlistadesimbolos
  % ---
  
  % ---
  % inserir o sumario
  % ---
  \pdfbookmark[0]{\contentsname}{toc}
  \tableofcontents*
  \cleardoublepage
  
}%hidelinks
% ---
